\scanpage{217}
\mspage{166}

(dans $\mathrm{Cat}_{/\unsure{X|x}}$, ou dans Cat, cela revient au $=$), et pour ceci j'applique le cor. 2 de la prop. 15 (p. 161), compte tenu que $X/q$ est un iso. local, et $X/i'$ (étant un iso. local mono) une immersion ouverte. Il suffit de prouver que, \deleted{\unsure{dans}} $T'$ et $T$ étant les \deleted{\unsure{fermés}} sous-catégories fermées complémentaires des ouverts $X/G'$ de $X/F'$ et $X/G$ de $X/F$, que $T\unsure{-|\longrightarrow}T'$ est iso.

La vérification est immédiate \unsure{;} $T$ est la cat. des \inserted{flèches} $x\longrightarrow G$ \deleted{qui} \inserted{dans $X^\wedge$} \inserted{(de source $x$ dans $X$)} qui ne se factorisent pas par $G$, $T'$ celle des flèches $x\longrightarrow G'$ qui ne se factorisent pas par $F'$,

\diagram{\begin{tikzcd}[ampersand replacement=\&,row sep=small,column sep=small]
F\arrow[r,hook,no head]\arrow[d] \& G\arrow[d]\\
F'\arrow[r,hook,no head]\& G'
\end{tikzcd}}

le foncteur $T\mathrel{\unsure{-|\longrightarrow}}T'$ étant le foncteur composition de $x\longrightarrow G$ avec $G\xrightarrow{q}G'$.

\deleted{\unsure{Or}} (NB le composé ne se factorise pas par $F'$, car autrement, à cause du caractère cartésien de $Q$, $x\longrightarrow G$ se factoriserait par $F$). Mais pour $x$ donné, la donnée de $\deleted{G(x)}$ $x\longrightarrow G$ ne se factorisant pas par $F$ est celle d'un élément $\deleted{\unsure{G|\in}}$
% archive sheet number: 217
